% =====================================================================
% overleaf/04_metrics_and_substrate.tex — LaTeX rendering of en/04_metrics_and_substrate.md v0.18 (17 September 2026)
% Dernière mise à jour : 2026-09-17
% Chapter 4 of the paper, a \section of its own, to be inserted right after overleaf/03_architecture.tex.
% This section DEFINES \label{sec:metrics}, which overleaf/03_architecture.tex already references
%   forward and which resolved to ?? until this file existed. Keep the label spelling exactly.
% Cross-references use \label / \ref, so no number is hard-coded.
% Forward reference: sec:results (chapter 6) is NOT defined yet, that chapter having no LaTeX
%   rendering so far. It resolves to ?? until 06 is written; the label is reserved with this name.
% Citations: Hörl & Balać 2021 (eqasim), Ke et al. 2017 (LightGBM), Meister et al. 2024 (verbalised
%   distributions). All three are already in references.bib — this file adds no entry.
%   NOT cited here: Martín-Baos et al. 2023. Two claims were attributed to it and both were wrong —
%   the paper finds the reverse illustration (the best DISAGGREGATE models lose on aggregate shares),
%   and it does not evaluate kernel logistic regression at all (MNL, SVM, RF, NN, XGBoost only; KLR
%   appears twice, in its related work, crediting Martín-Baos et al. 2020 and 2021). The claim on the
%   inversion of orderings is kept WITHOUT a source; the claim that KLR is the best accuracy /
%   behavioural-plausibility compromise is dropped, no source being identified for it (2026-09-16).
% Display equations: three, unnumbered (\[ ... \]), amsmath being loaded by aamas.cls.
%   They are definitions the text names, not results the text points back to.
% Tables: two, both \begin{table}[t] with the caption ABOVE, as acmart requires. booktabs is loaded
%   by aamas.cls. Table 1 fits one column; table 2 (demographic margins, 13 rows) is set as table*
%   across both columns — thirteen rows in one column would break the page.
% The survey footnote of the Markdown becomes a \footnote on the first mention of the microdata:
%   the ProGEDO deposit doi, required by agreement lil-1750.
% Comments "% source:" give the repository file a plain-text figure is recomputed from.
% Derived file: regenerate after every validated pass of en/04_metrics_and_substrate.md.
% =====================================================================


\section{Metrics and evaluation substrate}
\label{sec:metrics}

\subsection{The evaluation substrate}
\label{sec:substrate}

The measurements in this paper rest, for the most part, on one \textbf{test set}: a sealed cohort of 1,000 people, their trips on the evaluated day, and the itineraries the routers returned for each of them. It is on this set that most of the experiments of Sections~\ref{sec:prompt} and~\ref{sec:results} run, agents and tabular references alike. The few measurements carried out on a smaller set say so where they are cited. % source: data/jeux/population_1000_AAMAS_v6_20260316_EN/MANIFEST.yaml

The 1,000 people fall into 499 whole households, the household being the sampling unit, drawn from a pool of 11,329 people produced by the eqasim synthesis chain \citep{horl2021eqasim} from the census, tax income records and the household travel survey. The perimeter is the survey's: the 453 municipalities of the Toulouse living area. % source: data/population/population_1000_AAMAS_v6/MANIFEST.yaml, selection.menages_retenus.n = 499, selection.vivier.n = 11329

Each person carries a weight of one. The mobility of the cohort is that of its activity chains, which are cyclic: the day closes on the return home, so that a chain of $k$ activities carries $k$ trips. It produces 3,299 trips on the evaluated day, that is 3.30 per person and 3.69 per mobile person, with 10.6\,\% immobile; the survey gives 3.53, 3.95 and 10.6\,\%. % source: MANIFEST.yaml, controle.menages_et_mobilite and reference_enquete

Thirteen margins are controlled, each with its source: a page of the published report, or a recomputation on the survey microdata\footnote{Enquête Ménages Déplacements, Toulouse / Grande agglomération toulousaine, 2023. CEREMA and Syndicat mixte des transports en commun de l'agglomération toulousaine (producers), PROGEDO-ADISP (distributor), \texttt{doi:10.13144/lil-0933}. Used under agreement \texttt{lil-1750}: no microdata are redistributed.} (Appendix~\ref{app:margins}) where the report does not publish the margin.

\begin{table*}[t]
  \caption{Demographic control of the sealed cohort. Thirteen margins, each against the source that publishes it; the maximum absolute gap is given in percentage points. The test is an equivalence test margin by margin, with an indifference margin of 1 point announced before the measurement.}
  \label{tab:margins}
  \begin{tabular}{lllrl}
    \toprule
    Margin & Base & Source & Max gap (pt) & Verdict \\
    \midrule
    Age class (6 classes)                    & person    & published report & 0.50 & conforming \\
    Occupation                               & person    & published report & 0.50 & conforming \\
    Five-year age (15 classes)               & person    & microdata        & 0.29 & conforming \\
    Car ownership                            & person    & microdata        & 0.09 & conforming \\
    Car ownership                            & household & published report & 0.06 & conforming \\
    Residential ring                         & person    & published report & 0.06 & conforming \\
    Ring $\times$ car ownership              & person    & microdata        & 0.05 & conforming \\
    Household size                           & person    & microdata        & 0.06 & conforming \\
    Driving licence (18 and over)            & person    & microdata        & 0.01 & conforming \\
    Public transport season ticket           & person    & microdata        & 0.05 & conforming \\
    Dwelling type                            & person    & microdata        & 0.04 & conforming \\
    Gender                                   & person    & microdata        & 0.02 & conforming \\
    People with no trip on the previous day  & person    & microdata        & 0.05 & conforming \\
    \bottomrule
  \end{tabular}
  \Description{Table of the thirteen demographic margins controlled on the sealed cohort of 1,000 people. Each row gives the margin, the base on which it is computed (person or household), the source it is checked against (the published report or a recomputation on the survey microdata), the maximum absolute gap in percentage points, and the verdict. The gaps range from 0.01 to 0.50 point and all thirteen margins are judged conforming.}
\end{table*}

The test is an equivalence test margin by margin, with an indifference margin of 1 point announced before the measurement, completed by the maximum absolute gap and a Cramér's $V$ as effect size. % source: data/population/population_1000_AAMAS_v6/CONTROLE.md, sealed 2026-09-14

Trips are not independent: those made by the same person share their activity chain, their equipment and their home. Confidence intervals are therefore obtained by cluster resampling at the person level. The resolution reached is $\pm 1.3$ composite point per arm, and most of that noise comes from the least populated strata: the distance class beyond 50\,km carries only 17 decisions.
% source: docs/traces/2026-09-15_09-20_ticket080_idee_directrice_chapitre6/rapport.md § 2.1,
%   867 people drawn with replacement, 100 replicates, official scorer; paired differences at
%   200 replicates (§ 2.2). Measured before the rescoring of 2026-09-16: the composites there
%   run from 4.40 to 8.95 where the table of § 4.4 carries 3.60 to 4.09. The sizing note
%   computes on its side a design effect at household level (DEFF 2.1); the cluster used at
%   measurement is the person.

\subsection{What is measured}
\label{sec:quantities}

The population of agents chooses among the itineraries offered for each trip. Those choices are evaluated at two scales. At the \textbf{aggregate scale}, the question is fidelity: is the modal split produced that of the territory? At the \textbf{unit scale}, the question is agreement decision by decision: did this particular trip receive the mode the person declared?

\paragraph{The modal split.} Four modes are scored: car, public transport, walking, cycling. The reference is that of the EMC\textsuperscript{2} 2023 survey. It is recomputed for each stratum (age class, occupation, gender, trip purpose, distance class) and those readings are then aggregated into a composite score. % source: docs/arch/score-synthesis.md § 2 ; prompt_calibration/calibration/metrics.py

The \textbf{L1 error}, the sum of absolute share gaps, is expressed in percentage points and reads without conversion:
\[
  L_1 = \sum_{m} \left| \hat{P}_m - P_m^{\text{EMC}^2} \right|
\]
The \textbf{Jensen-Shannon divergence} is used on the nominal axes (gender, occupation, purpose):
\[
  \mathrm{JSD}(P \parallel Q) = \tfrac{1}{2}\,\mathrm{KL}(P \parallel M) + \tfrac{1}{2}\,\mathrm{KL}(Q \parallel M),
  \qquad M = \tfrac{1}{2}(P + Q)
\]
The \textbf{optimal transport distance} (Earth Mover's Distance, or Wasserstein-1 on a line) is used on the ordinal axes (age class, distance class). For $K$ ordered classes and $F$, $\hat{F}$ the cumulative distribution functions:
\[
  \mathrm{EMD}(\hat{P}, P) = \sum_{k=1}^{K-1} \left| \hat{F}_k - F_k \right|
\]
These readings aggregate into two composites. Each adds to the global reading the mean of each stratum, weighted 0.5 for age, occupation and purpose, 0.3 for gender and distance:
\[
  \mathcal{C}_{L_1} = L_1^{\text{global}} + \sum_{d} w_d\,\overline{L_1}^{\,d},
  \qquad
  \mathcal{C}_{\text{EMD--JSD}} = \mathrm{JSD}^{\text{global}}
  + \sum_{d\ \text{nominal}} w_d\,\overline{\mathrm{JSD}}^{\,d}
  + \sum_{d\ \text{ordinal}} w_d\,\overline{\mathrm{EMD}}^{\,d}
\]
where the bar is the unweighted mean over the categories of the stratum, categories of fewer than five personas being left out.
% source: prompt_calibration/calibration/metrics.py, WEIGHTS (global 1.0, age 0.5,
%   occupation 0.5, motif 0.5, genre 0.3, distance 0.3) and STRATUM_MIN_PERSONAS = 5. Two
%   historical terms, absent_penalty and length_penalty, are still computed and stored at zero
%   weight; the composite being linear, that removal re-applies exactly (rescale_composite).
%   Check on exp_lgbm_c_nosim of 2026-09-16: 9.491 + 0.5 (17.938 + 18.772 + 13.521)
%   + 0.3 (8.483 + 20.975) = 43.444.

\paragraph{Unit-level accuracy.} Three quantities are published at this scale: accuracy, cross-entropy, and recall and precision per mode.

\paragraph{The scored quantity.} A distribution over the four modes summarises in two ways, and both give a modal share: summing the probability assigned to each mode, or counting only the most probable mode. We keep the first, the \textbf{mass}. We ask the model to verbalise its distribution rather than to sample it by repetition (\citet{meister2024benchmarking} show that verbalised distributions align better than sampled ones), and the decision actually played in the simulation is a draw from that distribution. Keeping the most probable mode would flatten individual variability: every persona of the same profile would set off in the same mode, and the population would lose exactly what it is asked to reproduce. The mass is therefore the scored quantity, for every decision-maker compared, and it is what compares with the mass of a tabular model.

None of these quantities is sufficient on its own: the ordering of mode choice models reverses depending on the family of indicators, and a model can lead on aggregate shares while losing on disaggregate accuracy. That is why the quantities above are published together and read axis by axis, never summarised into a ranking.

\subsection{The three rules of the evaluation contract}
\label{sec:contract}

The first contribution announced in Section~\ref{sec:intro} comes down to three rules. Each closes one way in which a comparison between a generative agent and a tabular model can be biased.

\paragraph{Rule 1 — a common set of inputs.} The agent and the four tabular references receive the same 21 variables describing the person and the trip to be decided: twelve for the person and their household, three for the trip itself (purpose at origin, purpose at destination, departure hour), six for the geometry (distance between origin and destination, whether or not they lie in the same zone, distance to the centre from the origin and then from the destination, population density at each). The complete list is in Appendix~\ref{app:variables}. % source: scripts/progedo_logit/feature_spec.json

The agent additionally receives three things no tabular model can receive: the schedule of the trips it has left before returning home, the weather of the coming time slots, and the itineraries themselves, step by step, where the tabular model knows only the geometry of the trip and sees the offer only through the renormalisation of rule 3. % source: mobility_llm/src/mobility_llm/persona.py, fields agenda, day_outlook and trajectories
What is equalised is the 21 variables, not the information: the agent receives more, and that surplus is declared rather than corrected.

Three asymmetries frame the comparison, and they do not pull the same way. On exposure, the tabular references have read 39,203 survey trips and the agent has read none: that is what makes them high references rather than competitors. % source: scripts/progedo_logit/mode_choice_policy_metrics.json, training.n_fit 31,279 + training.n_valid 7,924.
% The sum is the right count for an exposure claim: early stopping (50 rounds, best_iteration 1500) lets the
% validation set choose the number of iterations, so it shapes the model. Section 1.3 carries the same 39,203 since chapter 1 v0.22.
On horizon, the agent sees its day and the tabular references see only one trip. On the offer, finally, the agent reads the itineraries proposed to it where the tabular references know them only through the modes they make available.

\paragraph{Rule 2 — the probability mass is the scored quantity.} The agent does not designate an itinerary: it spreads a probability mass over the options presented to it, six at most. That mass compares with the mass of the tabular model. The decision played in the simulation is a draw from it (Section~\ref{sec:quantities}, and Section~\ref{sec:decision}).

\paragraph{Rule 3 — renormalisation on the offer.} The tabular model predicts over four classes without knowing what was offered; the agent chooses only among the itineraries the routers actually returned. Every tabular prediction is therefore restricted to the modes offered for that trip, then renormalised to 100\,\%. The probabilities are written before and after correction. % effect measured, not published: renormalisation removes 1.5 point from gradient boosting's car share without the chain constraint (58.8 % to 57.3 %, n = 3,111) and 9.0 points under it (52.7 % to 43.7 %, n = 2,502)

\subsection{The four tabular reference models}
\label{sec:references}

Four classical mode choice methods serve as references: the multinomial logit, gradient boosting \citep{ke2017lightgbm}, kernel logistic regression and the random forest. Each is fitted on the survey microdata, then played on the test set of Section~\ref{sec:substrate}, trip by trip, exactly as the agents are. They are estimated at strict parity: same data file, same household split, same raising weights, same encoding, same metrics from a shared module. They are subject to the vehicle-chain constraint of Section~\ref{sec:chain}, the very one the agents are subject to: a car that left in the morning is not available in the evening.

\begin{table}[t]
  \caption{The four reference methods on the sealed cohort, under the vehicle-chain constraint. Lower is better on all three axes. In bold, the best score on each axis: it is that score, and not a designated model, that the ablation of Section~\ref{sec:results} compares against.}
  \label{tab:references}
  \small
  \setlength{\tabcolsep}{4pt}
  \begin{tabular}{@{}lrrr@{}}
    \toprule
    Reference method & \multicolumn{1}{c}{L1 global}
                     & \multicolumn{1}{c}{Composite}
                     & \multicolumn{1}{c}{Composite} \\
                     & \multicolumn{1}{c}{shares}
                     & \multicolumn{1}{c}{L1}
                     & \multicolumn{1}{c}{EMD--JSD} \\
                     & \multicolumn{1}{c}{(pts)}
                     & \multicolumn{1}{c}{(cum. pts)}
                     & \\
    \midrule
    Gradient boosting          & 9.5          & 43.4          & 3.60 \\
    Multinomial logit          & 9.0          & 44.9          & 4.02 \\
    Kernel logistic regression & 6.7          & 41.1          & 3.61 \\
    Random forest              & \textbf{5.3} & \textbf{38.3} & 4.09 \\
    \bottomrule
  \end{tabular}
  \Description{Table of the four tabular reference methods scored on the sealed cohort under the vehicle-chain constraint, on three axes where a lower value is better. Gradient boosting scores 9.5 on the L1 of global shares, 43.4 on the composite L1 and 3.60 on the composite EMD-JSD. The multinomial logit scores 9.0, 44.9 and 4.02. Kernel logistic regression scores 6.7, 41.1 and 3.61. The random forest scores 5.3 and 38.3, the best values on the first two axes, and 4.09 on the third. The third axis separates gradient boosting from kernel logistic regression by two thousandths and designates neither. No method leads on all three axes.}
\end{table}
% sources: data/experiences/exp_{lgbm,mnl,rf,klr}_jtir_pop-1000_AAMAS_v6_jeu-20260316_EN_c_nosim, scores.json:
%   global.l1, composite.l1, composite.emd_jsd. The reading without the chain constraint is not published:
%   it does not compare with the agents, all of which are subject to the chain.
%   Target: 56.7 / 12.4 / 26.8 / 4.1 % (car / public transport / walking / cycling).

No method leads on all three axes. The random forest has the lowest error on global shares, 5.3 points, and the best composite L1. On the composite EMD--JSD, gradient boosting and kernel logistic regression stand at 3.604 and 3.605: two thousandths separate them, for a resolution of $\pm$1.3 point per arm, and the axis therefore designates neither. Gradient boosting ends last on global shares, at 9.5 points. Kernel logistic regression, non-linear like the booster and smooth like the logit, follows the random forest on the other two axes. % source: klr_model_metrics.json ; RBF kernel, Nyström approximation

The four methods have read 39,203 survey trips where the agent has read none: together they form the high reference. The ceiling then reads axis by axis, and on each is worth the best score reached by one of the four: 3.60 of composite EMD--JSD, reached by gradient boosting and kernel logistic regression alike, 38.3 of composite L1 and 5.3 points of error on global shares for the random forest. It is against that score that an agent compares.
